% !TEX root = ../../main.tex

\chapter{语义分割}

我们之前介绍的卷积神经网络，都是用于图像分类这一任务的，其基本架构如\cref{fig:cnn_for_classification}所示：网络通过卷积层和池化层逐步提取图像特征，并在最后通过全连接层和分类器进行预测，输出图像所属的类别。

\begin{figure}[H]
    \boxedimage[width=0.95\textwidth]{figures/part_02/p2_ch06/cnn_for_classification.pdf}
    \caption{用于图像分类的卷积神经网络。输入$H\times W\times 3$图像经多轮“卷积+池化”逐步提取低、中、高层特征，再展平送入全连接层，最终由Softmax输出各类别概率。}
    \label{fig:cnn_for_classification}
\end{figure}

但是，视觉任务并不只有分类。例如，我们不仅希望知道一张图片中“有一只猫”，还希望进一步找出猫在图像中的具体位置，将属于猫的像素从背景中分离出来，如\cref{fig:segment}所示。这样的任务就是\textbf{图像分割（Image Segmentation）}。根据任务的不同，图像分割又可以分为\textbf{语义分割（Semantic Segmentation）}和\textbf{实例分割（Instance Segmentation）}。例如，当一张图片中有两只猫时，语义分割只需要判断每个像素属于猫还是背景，将两只猫都标记为“猫”即可；实例分割则需要进一步区分这两只猫，分别标记为“猫1”和“猫2”。接下来，我们看看卷积神经网络如何用于语义分割这一任务。

\begin{figure}[H]
    \boxedimage[width=0.5\textwidth]{figures/part_02/p2_ch06/segment.pdf}
    \caption{语义分割任务。根据图像内容，将输入$H\times W\times 3$图像划分为语义不同的区域：本例中模型提取猫所在区域并与背景分离，猫轮廓内标为前景（红）、其余标为背景。}
    \label{fig:segment}
\end{figure}

\section{全卷积神经网络}

语义分割需要对图像中的每个像素进行预测，判断其属于猫、背景或其他类别。因此，网络最终输出的也不再是一个表示整张图像的类别向量，而应该是一张与输入图像具有对应空间位置的预测图。这样一来，传统的分类网络就遇到了一个直接的问题：如果仍然按照原来的方式将特征图展平，再通过全连接层和分类器输出类别，那么原有的空间信息就会被破坏，网络最终只能得到整张图像的分类结果，而无法对每个像素分别进行预测。那么，我们能不能在分类网络的基础上进行改造，使网络在保留空间信息的同时输出一张像素级的预测图呢？

\subsection{$1\times1$卷积}

首先，我们构造一个简单的卷积池化网络，如\cref{fig:conv_pool}所示。输入图像大小为$H\times W\times3$，经过四次“卷积+池化”后，特征图的空间尺寸逐步减小，而通道数逐步增加，最终得到一张低分辨率的深层特征图：
$$
H\times W\times3 \rightarrow \frac{H}{2}\times\frac{W}{2}\times C_1 \rightarrow \frac{H}{4}\times\frac{W}{4}\times C_2 \rightarrow \frac{H}{8}\times\frac{W}{8}\times C_3 \rightarrow \frac{H}{16}\times\frac{W}{16}\times C_4
$$
随着网络逐渐加深，特征图的空间尺寸不断减小，而通道数逐渐增加，使每个空间位置能够包含更加丰富的语义信息。

\begin{figure}[H]
    \boxedimage[width=0.95\textwidth]{figures/part_02/p2_ch06/conv_pool.pdf}
    \caption{卷积池化网络。输入$H\times W\times3$图像经过四次“卷积+池化”，空间尺寸依次缩小为原来的$1/2$、$1/4$、$1/8$和$1/16$，同时通道数逐步增加为$C_1$、$C_2$、$C_3$和$C_4$，得到具有较强语义信息的深层特征。}
    \label{fig:conv_pool}
\end{figure}

经过四次池化后，特征图的空间尺寸已经缩小到原图的$1/16$，如果继续下采样到$1/32$，空间分辨率会进一步降低，预测结果也会更加粗糙，因此我们在这里停止下采样。另外，语义分割虽然名为“分割”，但本质上仍然是逐像素分类，即判断每个像素属于哪一类别，例如猫、狗、背景等。因此我们仍然需要进行分类，只是不再使用传统的全连接层，因为全连接层会将特征图展平，破坏原有的空间信息。那么，如何在保留空间结构的同时完成分类，又不进一步改变特征图的空间尺寸呢？

我们可以使用$1\times1$卷积。它在每个空间位置上对通道进行线性组合，在保留空间位置的同时完成分类。如果需要预测$K$个类别，就使用$K$个$1\times1$卷积核，为每个位置分别输出$K$个类别的得分：
$$
\frac{H}{16}\times\frac{W}{16}\times C_4 \xrightarrow{1\times1\ \mathrm{Conv}} \frac{H}{16}\times\frac{W}{16}\times K
$$
其中，$K$表示需要预测的类别数量。例如，如果区分“猫、狗、背景”三个类别，则$K=3$。特征图中的每个空间位置都会得到3个数，分别表示该位置属于“猫”“狗”和“背景”的预测得分。假设某个位置经过$1\times1$卷积后得到
$$
(2.1,0.3,-0.7)
$$
那么这3个数就分别对应猫、狗和背景的预测得分，其中猫的得分最高，因此该位置最可能属于猫。

这样一来，网络就可以在每个空间位置上分别进行分类，并将不同类别的预测结果保存在特征图的通道维度中。最终得到的$\frac{H}{16}\times\frac{W}{16}\times K$特征图，可以看作$K$张分别对应不同类别的\textbf{得分图（Score Map）}。

但这还不是最终的分割结果，因为得分图的空间尺寸只有输入图像的$\frac{1}{16}$。经过前面的多次池化后，原始图像中的多个像素可能已经对应到同一个特征位置，因此当前的预测结果仍然比较粗糙，无法直接与输入图像中的每个像素对应。接下来，我们需要进一步将这张低分辨率的得分图恢复到输入图像的空间尺寸。

\subsection{上采样与转置卷积}

将特征图空间尺寸放大的过程称为\textbf{上采样（Upsampling）}。例如，我们可以将$\frac{H}{16}\times\frac{W}{16}$的特征图逐步放大到$\frac{H}{8}\times\frac{W}{8}$、$\frac{H}{4}\times\frac{W}{4}$和$\frac{H}{2}\times\frac{W}{2}$，最终恢复到$H\times W$。最直接的做法是根据已有特征值计算新的位置，例如最近邻插值和双线性插值。最近邻插值直接使用邻近位置的特征值，双线性插值则根据周围位置的特征值进行加权计算，二者都可以在不引入可学习参数的情况下扩大特征图的空间尺寸。不过，这些方法采用的都是预先设定好的固定规则，上采样过程本身无法根据数据进行调整。

对于语义分割而言，我们希望网络能够根据训练数据学习低分辨率特征应该如何映射到更高分辨率空间，而不仅仅是按照固定规则进行放大。因此，我们需要让上采样过程本身具有可学习的参数，让网络在恢复空间分辨率的同时，也能够学习如何生成更加合理的高分辨率特征。

我们可以使用\textbf{转置卷积（Transposed Convolution）}。为了直观理解这一操作，我们先考虑一个简单的例子。假设输入是一个单通道的$2\times2$特征图：
$$
X=\begin{bmatrix}
x_1 & x_2\\
x_3 & x_4
\end{bmatrix}
$$
现在使用一个大小为$2\times2$的转置卷积核，并将步长设置为$2$。这个卷积核包含4个可学习参数：
$$
K=\begin{bmatrix}
k_1 & k_2\\
k_3 & k_4
\end{bmatrix}
$$
转置卷积的计算可以理解为：输入特征图中的每一个位置，都通过这个可学习的卷积核向输出空间中的一个局部区域产生贡献。例如，位置$x_1$经过卷积核$K$后，会在输出中产生一个$2\times2$的局部响应：
$$
\begin{bmatrix}
k_1x_1 & k_2x_1\\
k_3x_1 & k_4x_1
\end{bmatrix}
$$
同样，$x_2$、$x_3$和$x_4$也会分别经过同一个卷积核，在输出空间中产生对应的局部响应。由于步长为$2$，这四个$2\times2$区域恰好不会发生重叠，可以直接拼接成一个$4\times4$的输出：
$$
\begin{bmatrix}
k_1x_1 & k_2x_1 & k_1x_2 & k_2x_2\\
k_3x_1 & k_4x_1 & k_3x_2 & k_4x_2\\
k_1x_3 & k_2x_3 & k_1x_4 & k_2x_4\\
k_3x_3 & k_4x_3 & k_3x_4 & k_4x_4
\end{bmatrix}
$$
可以看到，原来的$2\times2$特征图经过转置卷积后变成了$4\times4$特征图。这里并不是简单地复制或插值，而是每个输入位置都通过同一个可学习的$2\times2$卷积核生成一个局部响应，因此网络可以通过训练自动学习如何将低分辨率特征展开到更高分辨率空间。

\begin{figure}[H]
    \boxedimage[width=0.95\textwidth]{figures/part_02/p2_ch06/transpose_convolution.pdf}
    \caption{转置卷积的具体例子（输入$2\times 2$、核$2\times 2$）。每个输入元素（彩色格）与核$K$相乘生成同色的$2\times 2$响应块，再按步长决定的偏移放置并相加。上排步长1：响应块逐格错位、彼此重叠，重叠位置累加（如中心15），得$3\times 3$输出；下排步长2：响应块互不重叠、直接拼接，得$4\times 4$输出。}
    \label{fig:transpose_convolution}
\end{figure}

在实际网络中，转置卷积通常处理多通道特征，并通过设置合适的卷积核大小、步长和填充来控制输出尺寸。例如，可以将低分辨率特征图逐步从$\frac{H}{16}\times\frac{W}{16}$放大到$\frac{H}{8}\times\frac{W}{8}$、$\frac{H}{4}\times\frac{W}{4}$，最终恢复到$H\times W$。当卷积核大小大于步长时，不同输入位置产生的响应还会在输出空间中发生重叠，并在对应位置进行累加，因此转置卷积本质上是利用可学习的卷积核，将输入特征逐步映射到更高分辨率的空间，而不是简单地复制或插值。stride为$1$和$2$时转置卷积的具体计算过程如\cref{fig:transpose_convolution}所示。

\begin{orangeBox}{为什么叫“转置卷积”？}
转置卷积中的“转置”，指的是卷积所对应的线性变换进行转置。为了直观理解这一点，我们直接来看一个二维卷积的简单例子。假设输入特征图为$X$，卷积核为$K$，则：
$$
X=\begin{bmatrix}x_1 & x_2\\ x_3 & x_4\end{bmatrix}\qquad K=\begin{bmatrix}k_1 & k_2\\ k_3 & k_4\end{bmatrix}
$$
设置步长为$2$且不使用填充，此时整个$2\times2$特征图只进行一次卷积，得到一个输出值：
$$
y=k_1x_1+k_2x_2+k_3x_3+k_4x_4
$$
将输入和输出分别展开成向量后，这个二维卷积可以写成矩阵乘法的形式：
$$
\begin{bmatrix}y\end{bmatrix} =
\underbrace{
\begin{bmatrix} k_1 & k_2 & k_3 & k_4 \end{bmatrix}}_{C}
\begin{bmatrix}
x_1 \\ x_2 \\ x_3 \\ x_4
\end{bmatrix}
$$
这里的矩阵$C$表示整个卷积操作。现在，将这个矩阵进行转置，可以得到：
$$
C^T=\begin{bmatrix} k_1\\ k_2\\ k_3\\ k_4 \end{bmatrix}
$$
将原卷积得到的输出$y$作为输入，经过转置后的线性变换$C^T$后，可以得到：
$$
C^Ty=\begin{bmatrix} k_1y\\ k_2y\\ k_3y\\ k_4y \end{bmatrix}
$$
再把这4个值重新排列成二维特征图，就得到：
$$
\begin{bmatrix}
k_1y & k_2y\\
k_3y & k_4y
\end{bmatrix}
$$
可以看到，普通卷积中，多个输入位置共同决定一个输出位置；而将对应的线性变换转置后，则变成了一个输入位置向多个输出位置传播信息。这正是转置卷积的核心特征，因此这种操作被称为\textbf{转置卷积（Transposed Convolution）}。在实际应用中，可以通过设置卷积核大小、步长和填充来控制输出特征图的空间尺寸，从而实现上采样。无论输出尺寸如何变化，其数学本质始终是卷积线性算子的转置。
\end{orangeBox}

现在，我们可以使用步长为$2$的转置卷积逐步将特征图恢复到$\frac{H}{8}\times\frac{W}{8}$、$\frac{H}{4}\times\frac{W}{4}$和$\frac{H}{2}\times\frac{W}{2}$，最终得到$H\times W$的预测图。这样，网络就能够为输入图像中的每个像素输出对应的类别得分。

不过，需要注意的是，转置卷积虽然能够通过可学习的参数将低分辨率特征图逐步映射到更高分辨率空间，却无法恢复下采样过程中已经丢失的原始像素信息。也就是说，它解决的是空间分辨率恢复的问题，而不是原始信息恢复的问题。

\subsection{跳跃连接}

根据前面的思路，我们已经可以构造出如\cref{fig:conv_pool_transpose}所示的语义分割网络：输入图像经过多次卷积和池化得到低分辨率的深层特征，再通过$1\times1$卷积得到类别得分，最后使用多次转置卷积将得分图逐步恢复到原始空间尺寸。这个网络已经能够完成像素级预测，但转置卷积只能根据当前保留下来的低分辨率特征生成更高分辨率的特征，而下采样过程中丢失的空间细节并不会因此恢复，因此最终得到的分割结果仍然比较粗糙。

\begin{figure}[H]
    \boxedimage[width=0.95\textwidth]{figures/part_02/p2_ch06/conv_pool_transpose.pdf}
    \caption{语义分割网络。下采样段经四层“卷积+池化”将分辨率逐次降至$1/16$（$\tfrac{H}{16}\times\tfrac{W}{16}\times C_4$），再由$1\times 1$卷积调整为$K$个通道；上采样段经四层转置卷积将分辨率逐次翻倍，恢复至$H\times W$，输出逐像素的$K$类预测。}
    \label{fig:conv_pool_transpose}
\end{figure}

既然深层特征在下采样过程中已经丢失了一部分空间信息，那么我们可以把前面那些仍然保留着空间细节的特征直接接到后面。这正是\textbf{跳跃连接（Skip Connection）}的基本思想。类似的思想我们在ResNet中已经见过，不过两者解决的问题并不完全相同：ResNet主要利用跳跃连接改善深层网络的优化，而语义分割中的跳跃连接则用于融合不同层次的特征，让深层特征提供丰富的语义信息，同时利用浅层特征补充更加精细的空间细节。加入跳跃连接后，整个网络如\cref{fig:full_convolution_network}所示，从输入到输出都不再需要全连接层，而是始终以特征图的形式进行计算，因此这种网络被称为\textbf{全卷积神经网络（Fully Convolutional Network）}\footnote{全卷积神经网络由Jonathan Long等人于2014年提出（arXiv预印本于2014年发布），论文为\href{https://openaccess.thecvf.com/content_cvpr_2015/papers/Long_Fully_Convolutional_Networks_2015_CVPR_paper.pdf}{\textit{Fully Convolutional Networks for Semantic Segmentation}}。原始论文为了复用VGG-16等分类网络在ImageNet上预训练的全连接层权重，采用了一种工程上的等价替换技巧：将全连接层改写成卷积核大小等于输入特征图空间尺寸的卷积层，从而在保留预训练参数的同时使网络能够输出空间得分图。本节介绍的是现代简化的FCN讲法，直接去掉全连接层并用$1\times1$卷积输出得分图，不再涉及原始论文中为复用VGG权重而做的具体实现细节。}。

\begin{figure}[H]
    \boxedimage[width=0.95\textwidth]{figures/part_02/p2_ch06/full_convolution_network.pdf}
    \caption{全卷积网络。相比\cref{fig:conv_pool_transpose}增加了跳跃连接（虚线）：将下采样段中保留空间细节的浅层特征（$1/2$、$1/4$、$1/8$分辨率）接入上采样段的同分辨率特征图进行融合，补偿下采样丢失的空间信息；网络从输入到输出不含全连接层，始终以特征图形式计算。}
    \label{fig:full_convolution_network}
\end{figure}

可以看到，越靠近网络后面的特征图，空间分辨率越低，但经过更多层卷积后具有更强的语义信息；越靠近网络前面的特征图，虽然语义信息相对简单，却保留了更多的位置、边缘和纹理信息。因此，我们可以在上采样过程中逐步引入浅层特征：先将$1/16$分辨率的深层特征上采样到$1/8$，与对应的浅层特征融合；再继续上采样到$1/4$，融合更浅层的特征，最后恢复到原始图像的空间尺寸。而为了使不同层的特征能够进行融合，可以先通过$1\times1$卷积调整浅层特征的通道数，使其与当前特征保持一致，再进行逐元素相加。

这样一来，网络就形成了一条完整的处理流程：先通过卷积和池化逐步提取深层特征，再利用$1\times1$卷积得到低分辨率的类别得分，随后通过转置卷积逐步恢复空间分辨率，并在上采样过程中不断融合浅层特征。深层特征负责判断“是什么”，浅层特征则帮助确定“在哪里”，二者结合后，网络就能够得到更加精细的像素级分割结果。

\section{U-Net}

\cref{fig:full_convolution_network}的全卷积网络实际上已经具有\textbf{编码器—解码器（Encoder-Decoder）}的基本思想：前面的卷积和池化不断降低特征图的空间分辨率，并提取越来越丰富的语义信息，这一部分可以看作\textbf{编码器（Encoder）}；后面的转置卷积则逐步恢复空间分辨率，将深层特征重新映射到原始图像空间，这一部分可以看作\textbf{解码器（Decoder）}。既然这种“先压缩、再恢复”的结构适合语义分割，那么能不能从一开始就围绕Encoder-Decoder结构专门设计一个网络，并进一步加强编码器与解码器对应层之间的特征连接呢？这就是\textbf{U-Net}的基本思想\footnote{U-Net由Olaf Ronneberger等人于2015年提出，论文为\href{https://arxiv.org/abs/1505.04597}{\textit{U-Net: Convolutional Networks for Biomedical Image Segmentation}}。}。

\begin{figure}[H]
    \boxedimage[width=0.95\textwidth]{figures/part_02/p2_ch06/unet.pdf}
    \caption{U-Net}
    \label{fig:unet}
\end{figure}

与全卷积网络直接在原有分类网络的基础上进行改造不同，U-Net从一开始就针对像素级预测设计了完整的Encoder-Decoder结构。如\cref{fig:unet}所示，左侧的编码器通过卷积和池化逐步降低空间分辨率，提取越来越丰富的语义特征；右侧的解码器则通过上采样逐步恢复空间分辨率，并在恢复过程中融合编码器对应层的浅层特征。由于整个网络从左到右逐步压缩，再从右向左逐步恢复，整体结构呈现出类似字母“U”的形状，因此被称为U-Net。

为了更清楚地理解U-Net的结构，我们按照网络的数据流逐层来看。原始U-Net使用大小为$572\times572\times1$的输入图像，并采用不填充的$3\times3$卷积，因此每经过一次卷积，特征图的空间尺寸都会缩小$2$个像素。整个网络可以概括为：

\begin{enumerate}
    \item \textbf{输入层}：接收$572\times572\times1$的输入图像。
    \item \textbf{编码器第1层}：连续使用两个$3\times3$卷积，每次卷积后经过ReLU激活，特征图尺寸由$572\times572$变为$568\times568$，通道数由$1$增加到$64$；随后使用$2\times2$最大池化，步长为$2$，空间尺寸缩小为$284\times284$。
    \item \textbf{编码器第2层}：连续使用两个$3\times3$卷积，特征图尺寸由$284\times284$变为$280\times280$，通道数增加到$128$；随后使用$2\times2$最大池化，得到$140\times140\times128$。
    \item \textbf{编码器第3层}：连续使用两个$3\times3$卷积，特征图尺寸由$140\times140$变为$136\times136$，通道数增加到$256$；随后使用$2\times2$最大池化，得到$68\times68\times256$。
    \item \textbf{编码器第4层}：连续使用两个$3\times3$卷积，特征图尺寸由$68\times68$变为$64\times64$，通道数增加到$512$；随后使用$2\times2$最大池化，得到$32\times32\times512$。
    \item \textbf{瓶颈层}\footnote{在编码器-解码器结构中，编码器不断下采样，空间分辨率逐渐降低；解码器则逐步上采样，恢复分辨率。两者之间的中间层通常被称为\textbf{瓶颈层（Bottleneck）}，因为它的空间尺寸最小，就像瓶子的颈部一样。瓶颈层虽然分辨率低，但通道数多、语义信息强，是整个网络“先压缩、再恢复”的关键转折点。}：连续使用两个$3\times3$卷积，特征图尺寸由$32\times32$变为$28\times28$，通道数增加到$1024$。这里已经得到整个编码器最深层的特征表示。
    \item \textbf{解码器第1层}：使用$2\times2$转置卷积，将$28\times28\times1024$上采样为$56\times56\times512$，通道数降低为$512$。随后，将编码器第4层得到的$64\times64\times512$特征图裁剪为$56\times56\times512$，再与上采样结果沿通道维进行拼接，得到$56\times56\times1024$特征图；最后连续使用两个$3\times3$卷积，将其变为$52\times52\times512$。
    \item \textbf{解码器第2层}：使用$2\times2$转置卷积，将$52\times52\times512$上采样为$104\times104\times256$；将编码器第3层的$136\times136\times256$特征图裁剪为$104\times104\times256$，与上采样结果拼接后得到$104\times104\times512$；随后经过两个$3\times3$卷积，得到$100\times100\times256$特征图。
    \item \textbf{解码器第3层}：使用$2\times2$转置卷积，将$100\times100\times256$上采样为$200\times200\times128$；将编码器第2层的$280\times280\times128$特征图裁剪为$200\times200\times128$后进行拼接，得到$200\times200\times256$；随后经过两个$3\times3$卷积，得到$196\times196\times128$特征图。
    \item \textbf{解码器第4层}：使用$2\times2$转置卷积，将$196\times196\times128$上采样为$392\times392\times64$；将编码器第1层的$568\times568\times64$特征图裁剪为$392\times392\times64$后进行拼接，得到$392\times392\times128$；随后经过两个$3\times3$卷积，得到$388\times388\times64$特征图。
    \item \textbf{输出层}：使用$1\times1$卷积，将$388\times388\times64$特征映射为$388\times388\times2$。这里的$2$表示需要预测两个类别，因此每个像素最终都会得到两个类别的得分。
\end{enumerate}

需要注意的是，原始U-Net的输入大小为$572\times572$，而最终输出大小为$388\times388$，两者并不相同。这是因为原始U-Net中的$3\times3$卷积均采用\textbf{不填充（valid convolution）}的方式，每经过一次卷积，特征图的高和宽都会缩小$2$个像素。随着多次卷积不断进行，特征图的空间尺寸逐渐减小，因此最终输出的分割结果也小于输入图像。与此同时，这种设计也使得编码器和解码器中对应层的特征尺寸不同，因此进行跳跃连接时，需要先对编码器中的特征图进行裁剪，再与解码器的特征进行拼接。

可以看出，U-Net与前面全卷积网络中的跳跃连接有所不同：全卷积网络主要通过逐元素相加融合不同层的预测结果，而U-Net则直接将编码器和解码器的特征沿通道维进行\textbf{拼接（Concatenation）}，再通过后续卷积进一步学习如何利用这两部分信息。整体上，U-Net采用典型的Encoder-Decoder结构：编码器逐步降低空间分辨率并提取越来越丰富的语义特征，解码器逐步恢复空间分辨率，而跳跃连接则将编码器中不同尺度的高分辨率特征传递给解码器，使网络能够同时利用深层的语义信息和浅层的空间细节，从而得到更加精细的像素级预测。

\section{ASPP}

还有一个问题：卷积的感受野究竟应该有多大？我们一直使用$3\times3$卷积，它能够有效提取局部特征，但对于不同大小的目标，这样的感受野未必都适用。对于较小的目标，较小的感受野可能已经能够提供足够的信息；而对于汽车、道路、天空等较大的目标，仅观察局部区域可能难以准确理解其整体语义。因此，我们希望网络能够在不同尺度上观察周围的上下文信息，从而更好地适应不同大小的目标。

既然需要同时获取不同尺度的上下文信息，我们可以参考前面介绍过的Inception模块，让多个不同尺度的卷积分支并行工作。例如，使用$3\times3$、$5\times5$和$7\times7$卷积，可以得到不同大小的感受野。不过，较大的卷积核会带来更多的参数和更高的计算量，而我们更希望在尽量不增加参数量的情况下扩大感受野。我们可以使用前面介绍过的空洞卷积来替代大卷积核。空洞卷积通过在卷积核元素之间引入间隔，可以在不增加卷积核参数数量的情况下扩大感受野。因此，我们可以将$5\times5$和$7\times7$卷积替换为具有不同空洞率的$3\times3$空洞卷积分支，在保持较少参数和计算量的同时获得不同大小的感受野。这样既保留了多尺度卷积的思路，又避免了直接使用大卷积核带来的额外参数和计算开销。不同空洞率的卷积分支具有不同的感受野，使同一个空间位置能够同时观察不同范围的上下文：小空洞率更关注局部信息，大空洞率则能够捕获更大范围的语义信息。将这些不同尺度的特征结合起来，网络就能够同时利用局部、中等范围和大范围的上下文信息。

不过，即使使用较大的空洞率，这些分支看到的仍然主要是局部区域，缺少对整张图像的全局理解。对于道路、天空等具有明显整体语义的目标，仅依靠局部上下文可能仍然不够，因此我们还需要增加一个图像级特征分支，直接提取整张特征图的全局信息。我们可以使用全局平均池化（Global Average Pooling，GAP）：它对每个通道的所有空间位置进行平均，将$\frac{H}{r}\times\frac{W}{r}\times C$的特征图压缩为$1\times1\times C$，从而得到概括整张特征图的全局表示；再经过$1\times1$卷积并上采样回原空间尺寸，就可以将全局语义信息与不同尺度的局部特征结合起来。这样一来，网络不仅能够观察不同大小的局部区域，还能够利用整张图像提供的全局上下文信息。

\begin{figure}[H]
    \boxedimage[width=0.95\textwidth]{figures/part_02/p2_ch06/aspp.pdf}
    \caption{空洞空间金字塔池化（ASPP）。输入特征图并行经过五条分支：$1\times 1$卷积、空洞率$r=1,2,3$的$3\times 3$空洞卷积，以及全局平均池化（将整图压缩为$1\times 1$全局表示再上采样回原尺寸）；各分支分别观察不同范围的上下文信息，输出沿通道维拼接后经$1\times 1$卷积（BN+ReLU）融合，得到最终输出。}
    \label{fig:aspp}
\end{figure}

这样的结构就是\textbf{空洞空间金字塔池化（Atrous Spatial Pyramid Pooling，ASPP）}\footnote{空洞空间金字塔池化（ASPP）由Liang-Chieh Chen等人于2016年在DeepLabv2中提出，论文为\href{https://arxiv.org/abs/1606.00915}{\textit{DeepLab: Semantic Image Segmentation with Deep Convolutional Nets, Atrous Convolution, and Fully Connected CRFs}}。}，如\cref{fig:aspp}所示。ASPP的核心思想其实并不复杂：不是让一个卷积核决定“应该看多大范围”，而是让多个分支同时从不同尺度观察同一个位置。这样，小目标可以利用较小范围的局部信息，大目标可以利用更大范围的上下文信息，而图像级分支则进一步提供全局语义信息，将不同尺度的信息融合起来后，网络就能够获得更加丰富的多尺度特征表示。

\section{DeepLab}

我们已经逐步解决了语义分割中的几个关键问题：全卷积网络让卷积神经网络能够保留空间结构并进行逐像素预测，跳跃连接能够补充下采样过程中丢失的空间细节，Encoder-Decoder结构则让网络能够在提取深层语义信息的同时逐步恢复空间分辨率，而ASPP进一步让网络能够从不同尺度观察目标周围的上下文信息。现在，把这些思路结合起来，就可以得到一个更加完整的语义分割网络，这就是\textbf{DeepLabV3+}\footnote{DeepLab系列由Liang-Chieh Chen等人逐步发展：DeepLabv1（2015）引入空洞卷积和全连接条件随机场（CRF）后处理；DeepLabv2（2017）提出空洞空间金字塔池化（ASPP）；DeepLabv3（2017）改进ASPP并移除CRF，直接上采样输出；DeepLabv3+（2018）在此基础上引入编码器—解码器结构，融合浅层特征。本文介绍的DeepLabv3+论文为\href{https://arxiv.org/abs/1802.02611}{\textit{Encoder-Decoder with Atrous Separable Convolution for Semantic Image Segmentation}}。}。

\begin{figure}[H]
    \boxedimage[width=0.95\textwidth]{figures/part_02/p2_ch06/deeplab.pdf}
    \caption{DeepLabV3+整体架构。Encoder：骨干网络提取浅层特征（$\tfrac{H}{4}\times\tfrac{W}{4}\times 256$）与深层特征（$\tfrac{H}{16}\times\tfrac{W}{16}\times 2048$），深层特征经ASPP模块获得多尺度上下文信息。Decoder：ASPP输出先$\times 4$上采样，与经$1\times 1$卷积压缩至48通道的浅层特征沿通道拼接，再经$3\times 3$卷积精炼、$\times 4$上采样，得到最终逐像素预测（$H\times W\times K$）。}
    \label{fig:deeplab}
\end{figure}

如\cref{fig:deeplab}所示，DeepLabV3+整体采用Encoder-Decoder结构。输入图像首先经过DCNN（Deep Convolutional Neural Network）网络提取特征，同时保留不同层次的特征：高层特征的空间尺寸为$\frac{H}{16}\times\frac{W}{16}$、通道数为$2048$，具有较强的语义信息；浅层特征的空间尺寸为$\frac{H}{4}\times\frac{W}{4}$、通道数为$256$，保留了更加丰富的空间细节。高层特征首先进入ASPP，通过$1\times1$卷积、不同空洞率的$3\times3$空洞卷积以及图像级特征分支提取多尺度上下文信息，各分支融合后经过$1\times1$卷积得到$\frac{H}{16}\times\frac{W}{16}\times256$的特征。随后将其上采样$4$倍恢复到$\frac{H}{4}\times\frac{W}{4}$，同时将浅层特征通过$1\times1$卷积压缩到$48$个通道，再与上采样后的高层特征沿通道维进行拼接，得到$\frac{H}{4}\times\frac{W}{4}\times304$的融合特征。融合后的特征经过$3\times3$卷积进一步细化，再上采样$4$倍恢复到$H\times W$，最终输出每个像素对应的$K$个类别预测。

\begin{quote}
    至此，我们已经介绍完了语义分割的基本方法。当然，卷积神经网络的应用远不止于此，例如目标检测、实例分割、姿态估计等任务中都可以看到它的身影。受篇幅限制，这些内容我们暂时不展开了。接下来，我们看看神经网络是如何理解和处理文本数据的。
\end{quote}
